% Development scaffold. Unresolved results remain explicit.
% Build with scripts/build_paper_v26.py and the pinned, unmodified ACL style.
\documentclass[11pt]{article}
\usepackage[review]{acl}
\usepackage{times}
\usepackage{latexsym}
\usepackage[T1]{fontenc}
\usepackage[utf8]{inputenc}
\usepackage{microtype}
\usepackage{booktabs}
\usepackage{amsmath}
\hypersetup{pdfauthor={},pdftitle={When Recency Is Not Currentness},pdfsubject={Development analysis scaffold}}
\newcommand{\pending}[1]{\textbf{[PENDING: #1]}}

\title{When Recency Is Not Currentness:\\Auditing Temporal Interpretation in Retrieval-Augmented QA}
\author{Anonymous ACL submission}

\begin{document}
\maketitle
\begin{abstract}
\textbf{Development scaffold; natural diagnostic pending.}
Temporal question answering requires matching evidence to the requested entity, scope, and time.
Selecting the newest timestamp does not establish that match.
We audit our earlier experimental package and identify several confounded measurements.
Among 35 current-holder questions, nine selected reference claims do not explicitly support the requested current role.
A global date baseline retains the requested subject--relation key for only five questions.
The reported aggregate improvement cannot reach $p<0.05$ under any compatible exact paired McNemar test.
These findings concern one package, with assistant-assisted source review.
We then specify a development diagnostic on 60 unchanged natural questions.
It separates interpretation correction from additional evidence, ordinary re-extraction, and draft anchoring.
The natural results remain unmeasured here.
We establish neither a new retrieval algorithm nor an answer-level advantage.
\end{abstract}

\section{Introduction}

An older document may describe a continuing state.
A newer document may describe a different population or a past event.
Before suppressing evidence, a temporal system must establish what the question and sources mean together.
This includes distinguishing an event's start from a document's publication and a claim's continuing validity.
These distinctions motivate our evaluation; they are not proposed as new temporal concepts.

Recent methods already address substantial parts of this problem.
MRAG separates question content from temporal constraints and ranks evidence using semantic and temporal signals \citep{zhang2025mrag}.
Re$^3$ combines temporal representations with conflict-aware recency filtering \citep{cao2026re3}.
Its recency rule assumes that the newest conflicting version supersedes earlier versions within credible corpora.
Chronos organizes evolving evidence into an event graph \citep{liu2026chronos}.
Adding temporal fields or a revision prompt therefore does not, alone, establish a new method.

We instead ask what an evaluation actually establishes before claiming better temporal reasoning.
Our earlier package compared lifecycle filtering with flat retrieval on current-holder questions.
Its apparent advantage could reflect reference construction, topical evidence loss, scoring, or temporal reasoning.
We disentangle those explanations through a retrospective self-audit.
We preserve every original reference and report all questions.

The audit motivates a second, prospective question.
When additional evidence arrives, does a model correct an asserted temporal interpretation beyond ordinary re-extraction?
We design a bounded natural-question diagnostic to test that premise before developing a specialized mechanism.
The completed contribution in this draft is the package audit.
The natural diagnostic and any broader scientific claim remain incomplete.

\section{Retrospective Package Audit}
\label{sec:audit}

\paragraph{Materials and procedure.}
The audited package contains claims, reference rules, retrieval code, aggregate results, and cached model outputs.
Its natural-prose comparison contains 35 current-holder questions.
None supplies a query date.
We inspect the reference construction and all 35 selected reference claims.
We additionally inspect complete same-key histories for nine flagged questions.
Source judgments are assistant-assisted annotations, not independent human gold.
Exact source positions and hashes support later review.
The audit makes no new model calls and checks no live-world officeholders.

\paragraph{Current support versus latest start.}
The reference rule selects the maximum timestamp within a curated subject--relation group.
Those timestamps encode tenure starts, rather than publication dates.
Table~\ref{tab:support} summarizes the selected claims' released text.
Nine claims do not explicitly assert the requested current role.
Six explicitly describe ended tenures.
The remaining cases involve past wording, another club, or an absent club name.
These observations establish internal task--evidence mismatches.
They do not identify the actual current holder.
Asking for the latest supplied start would define a different, narrower task.

\begin{table}[t]
\centering
\small
\begin{tabular}{lr}
\toprule
Selected reference claim's released text & Count \\
\midrule
Explicit requested current role & 26 \\
Explicit ended tenure & 6 \\
Most recent past role only & 1 \\
Different club named as current & 1 \\
Requested club absent & 1 \\
\midrule
All questions & 35 \\
\bottomrule
\end{tabular}
\caption{Assistant-assisted review of all selected reference claims. These are source-support categories, not external factuality labels.}
\label{tab:support}
\end{table}

\paragraph{A baseline without topical evidence.}
The date baseline sorts the entire corpus by timestamp before retaining five claims.
Exactly five claims have the maximum timestamp.
Every question therefore receives those same claims, although their order can differ.
Only five questions retain their requested subject--relation key.
Enumerating all possible orders recovers all 35 cached responses through exact prompt hashes.
The recovered score matches the original 5/35 count.
Thus, this baseline's failure cannot establish failure of subject-aware temporal selection.

\paragraph{Scoring and incomplete pairing.}
The original scorer accepts normalized reference containment or a shared name token with at least four characters.
Its stricter variant also rejects listed stale names.
These rules are not strict exact match or independently adjudicated identity correctness.
Full-name phrase matching and normalized equality are therefore reported separately in the audit artifact.
Their disagreements cannot automatically be called false positives.

The cache stores 4,995 unique prompt hashes, but neither prompts nor condition labels.
Bounded reconstruction recovers 139 compatible condition records from 175 possible cells.
These records use 130 distinct cache keys; shared prompts are not independent generations.
Another 29 cells remain unrecovered and seven remain condition-ambiguous.
Missing reconstruction does not establish missing generation.
Flat and lifecycle conditions also use different prompt wording.
The original comparison consequently does not isolate retrieval under an identical reader prompt.

\paragraph{What the aggregate improvement permits.}
The original report gives 30/35 lifecycle successes and 25/35 flat successes.
These margins permit six paired contingency tables.
Exact two-sided McNemar $p$-values range from $0.0625$ to $0.3018$ across those tables.
Even the most favorable pairing does not meet $0.05$ under this test.
This result does not establish equivalence or prove that temporal filtering cannot help.
The 19 reconstructed lifecycle--flat pairs contain two wins, zero losses, and 17 ties under the original scorer.
Their exact paired $p$-value is $0.5$.
This selected subset cannot estimate the missing pairs.

\paragraph{Simpler controls and timestamp precision.}
An existing variant removes older matched claims without comparing their values.
Its mask equals the original mask on five of seven inspected datasets, including the natural-prose set.
With supplied subject--relation keys, newest-value selection matches all 35 natural-prose references by construction.
That oracle-conditioned result is not natural retrieval accuracy.
It shows why key identification and temporal selection need separate measurements.

The original ISO parser also loses day information in 151 of 153 complete-date Wikidata records.
Its boundary condition fails before the ISO time separator and falls back to the year.
Eight groups collapse distinct dates and values into one timestamp.
Using supplied calendar dates changes two selected references.
Eighteen partial dates remain partial; we do not invent missing calendar fields.
The prose derivative shares these histories and is not an independent confirmation.
No original references are overwritten, and no corrected-date reader outputs are generated.

\section{A Natural Interpretation Diagnostic}
\label{sec:diagnostic}

The audit does not reveal whether interpretation revision offers useful headroom on natural questions.
We specify the following development study to answer that narrower question.
Protocol details below are prospective unless explicitly marked completed.

\paragraph{Questions and admission.}
TEMPO contains 1,730 naturally occurring Stack Exchange questions across 13 domains \citep{abdallah2026tempo}.
Its documents and assisted annotations support a natural-question diagnostic, rather than a source-revision replay.
Our frozen roster contains 60 original questions, selected through a recorded hash rule and domain round-robin ordering.
Six previously inspected questions are excluded before sampling.
Sampling uses identifiers and domains, without predictions or question contents.
All 60 remain in the accounting, including unanswerable, ineligible, or failed items.
This domain-balanced sample does not estimate the benchmark's natural domain mixture.

Question-only admission checks whether the original request has an objectively reviewable temporal target.
Merely mentioning a duration does not establish such a target.
Broad causal explanations and code review cannot silently become narrower date questions.
We record representation limits, missing anchors, and any dependence on inaccessible images.
These assistant annotations do not enter inference inputs.

\paragraph{Interpretations and evidence.}
An interpretation assigns entity, relation, scope, time, comparison, and replacement.
Unknown values differ from inapplicable fields.
Two apparently conflicting claims may coexist because their scopes or periods differ.
An older timestamp alone does not establish replacement.
This schema is a measurement instrument, not a novel temporal representation.

Each admitted item receives immutable packs $C_0\subseteq C_1$ before predictions.
The expanded pack cannot depend on an initial model error.
Every unit retains source identifiers, positions, and hashes.
Question dates, publication dates, and omitted premises remain unknown unless supplied.
If gold document identifiers define a diagnostic pool, we disclose that conditioning.
Such a pool cannot establish ordinary retrieval performance.
Benchmark answers and reference assignments never enter model inputs.
Missing evidence within $C_1$ cannot establish global unanswerability.

\paragraph{Matched controls.}
Initial extraction, $I_0$, receives $C_0$.
Ordinary re-extraction, $B_1$, receives $C_1$ and treats $I_0$ as an untrusted draft.
Candidate-preserving revision, $R_1$, receives identical evidence and the same draft.
It records retained, changed, retired, and newly proposed interpretations.
Fresh extraction, $F_1$, receives $C_1$ without the draft.
This control tests whether apparent revision difficulty instead reflects anchoring.
All arms share the model, decoding, output schema, candidate cap, and output budget.
The revision prompt is an exploratory control, not an algorithmic contribution.

Source-assisted references are written from $C_1$ before predictions and reviewed by a second role.
They specify acceptable joint assignments and supporting evidence sets.
Assistant-authored references remain distinct from independent human gold.
At most two active interpretations fit the proposed schema.
References requiring more alternatives remain recorded representation exclusions.

\paragraph{Outcomes and stopping rule.}
An initial error requires an asserted wrong assignment under the reviewed expanded evidence.
An unknown field becoming known measures evidence gain, not interpretation correction.
We report complete joint correctness, observed class recall, false elimination, corrections, and new harms.
Frozen aliases support strict scoring; blinded review audits apparent gains and losses.
Unresolved measurement disagreements cannot count as mechanism wins.
Malformed outputs count as failures without model-based repair.

The feasibility gate requires at least 12 source-resolvable initial errors among the 60 questions.
At least six must remain wrong after $B_1$ despite a complete reviewed correction from $C_1$.
We report the same residual against $F_1$.
If fresh extraction removes the residual, specialized revision loses its practical motivation here.
These thresholds are development decisions, not statistical evidence of superiority.
All calls, repeated evidence, tokens, and latency are charged.
The two revision arms include initial-extraction cost; fresh extraction retains its cheaper native cost.

\section{Natural Diagnostic Results: Pending}
\label{sec:pending}

\pending{N1: admission counts and independently reviewed source sufficiency across all 60 questions.}

\pending{N2: frozen runtime, exact evidence budgets, completed arm records, failures, and total costs.}

\pending{N3: initial errors, correction and harm counts, class recall, and residuals against both controls.}

These placeholders are not zero-valued observations.
No natural interpretation advantage, answer-level gain, or learned mechanism is established in this draft.
A later answer-level comparison also requires a separate reader-competence gate.
Passing the diagnostic would authorize further research, rather than establish publication readiness.

\section{Conclusion}

Our package audit separates current support from the latest supplied start date.
It exposes topical baseline failure, incomplete pairing, permissive scoring, and timestamp precision loss.
These findings weaken the earlier empirical claim without establishing a replacement method.
The natural diagnostic is designed to determine whether interpretation revision deserves further development.
Its results and external validation remain necessary before a broader analysis claim is justified.
\label{v26-main-end}

\clearpage
\section*{Limitations}

This is a development scaffold with explicit unfinished results.
The completed audit concerns our own experimental package.
It does not measure the prevalence of these problems across published systems or public benchmarks.
Source-support judgments use assistant assistance and require independent human adjudication before stronger claims.
We have not checked the actual current officeholders or reconstructed missing archived publication dates.

Cache reconstruction is bounded, and the original cache omits condition labels and complete runtime provenance.
The recovered subset is selected by reconstructability.
Original name scoring also cannot establish strict semantic correctness.
The aggregate paired test addresses statistical uncertainty under that original scoring rule.
It cannot repair task-definition or measurement problems.

The proposed natural study is development material, with no held-out confirmation.
Its compact schema cannot represent every natural question.
Exact aliases can reject valid paraphrases, and observed interpretation classes need not exhaust all valid alternatives.
Pack sufficiency does not certify corpus coverage.
Static evidence expansion does not test historical source replacement or adaptive retrieval quality.
No natural method win has been demonstrated.

\section*{Research Assistance and Data Handling}

AI assistants supported literature review, source annotations, code preparation, and manuscript drafting.
Their annotations are explicitly labelled and are not independent human judgments.
This scaffold contains no personal identity or project repository link.
An anonymous submission supplement remains to be assembled and reviewed separately.
Public project files do not themselves constitute an anonymous supplement.
Third-party source text remains subject to its original rights.
Hashes and offsets support auditing without assuming permission to redistribute complete captures.

\bibliography{interpretation_v26}
\end{document}
